\section{Euclidean cost and a complete Gaussian flip}
\label{sec:cost-reflections}

\subsection{The cost convention}

For a nonzero positive vector $v=(y,x)^{\mathsf T}$, let $\ell(v)$ be the
number of unit subtractions in the algorithm that repeatedly subtracts
the smaller coordinate from the larger and stops when the coordinates
are equal. This definition is unchanged by interchanging coordinates or
multiplying both by a positive integer. For a vector with nonzero signed
coordinates, we apply it to their absolute values.

Write
\begin{equation}
 M(j)=\begin{pmatrix}j&1\\1&0\end{pmatrix},\qquad
 M(q_1\cdots q_r)=M(q_1)\cdots M(q_r),\qquad
 e_1=\begin{pmatrix}1\\0\end{pmatrix}.
 \label{eq:digit-matrix}
\end{equation}
For a positive continued-fraction word $q$ with final digit at least two,
the first column of $M(q)$ is a primitive vector $(y,x)^{\mathsf T}$ with
$y>x>0$. The usual canonical finite expansion of $y/x$ is $[q_1;\ldots,q_r]$.
Division with remainder gives
\begin{equation}
 \ell(M(q)e_1)=q_1+\cdots+q_r-1.
 \label{eq:euclidean-cost}
\end{equation}
Indeed, the successive divisions make $q_i$ subtractions until the last
division, where the algorithm stops at equal coordinates and makes one
fewer. The final minus one is part of our convention throughout.

Let $\chi$ double each digit, so that $\chi(121)=112211$. Our regular source
language is
\begin{equation}
 q=\varepsilon\chi(w)2,\qquad
 \varepsilon\in\{\varnothing,1,2\},\quad w\in\{1,2\}^*.
 \label{eq:regular-source}
\end{equation}
An empty optional digit contributes zero to a digit sum. If $w=j_1\cdots j_n$,
then
\begin{equation}
 \ell(M(q)e_1)=\varepsilon+2\sum_{i=1}^n j_i+1.
 \label{eq:paired-cost}
\end{equation}
The language in \eqref{eq:regular-source} is larger than the genuine
Markoff source language. This enlargement makes the endpoint theorem
stronger, while genuine sources will later supply balance and coprime
Farey counts.

\subsection{The two reflection numerators}

The rational prime $89$ splits in the ring $\Z[i]$ of Gaussian integers.
Fix once and for all the Gaussian prime above it given by
\begin{equation}
 \pi=8+5i,\qquad \bar\pi=8-5i,\qquad \pi\bar\pi=89.
 \label{eq:gaussian-prime}
\end{equation}
The four associates $\pm\pi,\pm i\pi$ generate the same ideal, and
$\bar\pi$ generates the other prime above $89$. Every choice made below is
expressed through this one Gaussian integer, which makes the construction
rigid: the coordinate pair $(8,5)$ of $\pi$ is itself the primitive vector
of the source $q=1112$, identified in \cref{cor:genuine-strict} as the
unique genuine source of vanishing excess. The two integral matrices
\begin{equation}
 R_0=\begin{pmatrix}39&80\\80&-39\end{pmatrix},\qquad
 R_1=\begin{pmatrix}80&39\\39&-80\end{pmatrix}
 \label{eq:reflections}
\end{equation}
satisfy $R_\nu^{\mathsf T}R_\nu=89^2I$ and $\det R_\nu=-89^2$.
Under the identification $(y,x)^{\mathsf T}\leftrightarrow y+ix$, they act
as $z\mapsto\pi^2\bar z$ and $z\mapsto i\bar\pi^2\bar z$, respectively.

For a nonzero integral vector $u$, write $\cont{u}$ for the greatest common
divisor of its coordinates, taken positive. The next lemma explains why
the exact endpoint content, rather than divisibility alone, matters.

\begin{lemma}[Exact primitive endpoint]
\label{lem:primitive-flip}
Let $v=(y,x)^{\mathsf T}$ be primitive, with $x,y>0$, and let
$N=x^2+y^2$. Then $\ord_{89}(N)=1$ if and only if exactly one of $R_0v,R_1v$
has content $89$. For this branch, the vector $R_\nu v/89$ is primitive,
has nonzero coordinates and norm $N$. Modulo signs and coordinate order,
it represents the complete flip of the Gaussian prime pair above $89$.
\end{lemma}

\begin{proof}
The prime $89$ splits as $\pi\bar\pi$ in $\Z[i]$. Since $v$ is primitive,
$z=y+ix$ cannot be divisible by both $\pi$ and $\bar\pi$. If $\pi^e$ divides
$z$ exactly, with $e>0$, then the valuations of $\pi^2\bar z$ at
$(\pi,\bar\pi)$ are $(2,e)$. Its rational content at $89$ is therefore
$89^{\min(2,e)}$. The other branch has content prime to $89$. The roles
reverse when $z$ is divisible by $\bar\pi$. When neither divides $z$,
neither branch has content divisible by $89$.

No rational prime other than $89$ can divide a reflected content: the
matrix is invertible modulo every such prime. Since $N=z\bar z$ and only
one of $\pi,\bar\pi$ divides $z$, the exponent $e$ is exactly
$\ord_{89}(N)$. Thus content exactly $89$ is equivalent to $\ord_{89}(N)=1$. Division gives a primitive norm-$N$ vector.
If one coordinate vanished, $N$ would be a square, contradicting
$\ord_{89}(N)=1$.

For example, if $z=\pi t$, the compatible branch divided by $89$ is
$\pi\bar t$. Its conjugate is $\bar\pi t$, which exchanges precisely the
allocation of $\pi$ and $\bar\pi$. Units and conjugation account for the
irrelevant signs and coordinate order. This is the stated complete flip.
\end{proof}

For the branch selected by \cref{lem:primitive-flip}, define the completed
cost excess
\begin{equation}
 \Delta=\ell(R_\nu v/89)-\ell(v).
 \label{eq:completed-excess}
\end{equation}
The normalization by $89$ does not change the subtractive cost. It does,
however, decide whether the terminal vector is the primitive
representation being compared. If $\ord_{89}(N)\geqslant2$, the compatible
raw content is $89^2$ and the present theorem does not apply.
